\documentclass[a4paper,12pt]{article}
\usepackage[utf8]{inputenc}
\usepackage[T1]{fontenc}
\usepackage[margin=2.5cm]{geometry}
\usepackage{ctex} % 支持中文
\usepackage{booktabs} % 专业表格线
\usepackage{array} % 表格列格式
\usepackage{graphicx} % 图片支持
\usepackage{titlesec} % 标题格式
\usepackage{xcolor} % 颜色支持
\usepackage{enumitem} % 列表优化
\usepackage{float} % 浮动体控制

% 设置标题格式
\titleformat{\section}{\Large\bfseries\color{blue!40!black}}{}{0.5em}{}
\titleformat{\subsection}{\large\bfseries\color{gray!80!black}}{}{0.5em}{}

% 列表设置
\setlist[itemize]{leftmargin=*}

% 元数据
\title{\textbf{2024年度环境、社会及管治（ESG）报告节选}}
\author{公司董事会办公室}
\date{2024年3月}

\begin{document}

\maketitle

\section{1. 前言：战略定位与可持续发展承诺}

2024年是海南自贸港封关运作准备的关键之年。作为深度参与自贸港建设的大型基础设施及临空产业上市公司，本公司（以下简称“公司”）紧密围绕“一主一次，两支一货运”的区域机场群布局，积极履行企业社会责任。公司致力于通过高标准的治理架构、环境管理体系及社会价值创造，服务国家开放战略，并为股东及利益相关方创造长期价值。

\section{2. 环境范畴（Environmental）：低碳运营与资源管理}

公司严格遵守国家及地方法律法规，制定并执行环境管理策略。2024年，公司持续推进“双碳”目标落地，在能源管理、温室气体减排及绿色建筑认证方面取得实质性进展。

\subsection{2.1 排放与资源消耗绩效}

报告期内，公司完成了年度温室气体排放核查及能源审计工作。主要环境绩效指标如下表所示：

\begin{table}[H]
    \centering
    \caption{2024年度环境绩效指标表}
    \label{tab:env_kpi}
    \renewcommand{\arraystretch}{1.2}
    \begin{tabular}{p{8cm} p{3cm} p{3.5cm}}
        \toprule
        \textbf{环境指标} & \textbf{数值} & \textbf{单位} \\
        \midrule
        \multicolumn{3}{l}{\textbf{温室气体排放}} \\
        温室气体排放总量 & 51,518.71 & 吨CO$_2$当量 \\
        \quad 其中：范围1（直接排放） & 11,214.39 & 吨CO$_2$当量 \\
        \quad 其中：范围2（间接排放） & 40,304.32 & 吨CO$_2$当量 \\
        排放强度 & 0.12 & 吨CO$_2$当量/万元 \\
        \midrule
        \multicolumn{3}{l}{\textbf{能源消耗}} \\
        综合能源消耗量 & 13,538.22 & 吨标准煤 \\
        外购电力 & 7,511.05 & 万千瓦时 \\
        汽油消耗 & 19.10 & 万升 \\
        \midrule
        \multicolumn{3}{l}{\textbf{水资源与废弃物管理}} \\
        水资源消耗总量 & 145.94 & 万吨 \\
        中水回用量 & 13.61 & 万吨 \\
        无害废弃物产生量 & 6,230.58 & 吨 \\
        有害废弃物产生量 & 0.58 & 吨 \\
        废弃物合规处理率 & 100 & \% \\
        \bottomrule
    \end{tabular}
\end{table}

\subsection{2.2 绿色机场建设与减排措施}

公司积极推进技术节能与管理节能，重点实施以下减排项目：

\begin{itemize}
    \item \textbf{APU替代设施应用}：
    \begin{itemize}
        \item 在旗下核心门户机场全面推广辅助动力装置（APU）替代设施，旨在减少航空器地面运行期间的燃油消耗及噪音污染。
        \item \textbf{实施规模}：投入使用APU替代设备共计31套（近机位18套，远机位13套）。
        \item \textbf{环境效益}：全年实现减碳量 \textbf{25,382} 吨CO$_2$当量。
        \item \textbf{认证荣誉}：旗下核心门户机场获评“双碳机场”二星级称号。
    \end{itemize}
    
    \item \textbf{可再生能源利用}：
    \begin{itemize}
        \item \textbf{商业综合体分布式光伏}：旗下核心商业综合体光伏项目年均发电量 450 万千瓦时，占其总用电量的10\%，年减排 2,263.59 吨CO$_2$。
        \item \textbf{临空产业园光伏}：完成屋顶光伏装机容量 8.6MWp，预计年减排贡献 25.75 万吨。
    \end{itemize}

    \item \textbf{场内车辆电动化}：
    \begin{itemize}
        \item 核心机场新能源车辆保有量达 221 辆，占比提升至 43.6\%；配套新建充电桩约 100 个。
        \item 车辆电动化全年实现减碳 761 吨。
    \end{itemize}
\end{itemize}

\subsection{2.3 绿色建筑认证}

公司在项目开发建设中严格遵循绿色建筑标准：
\begin{itemize}
    \item 省会城市在建地标项目“海南中心”（双子塔）获得 \textbf{LEED CS 金级预认证}。
    \item 某新区F07项目通过绿色建筑一星级标准认证。
\end{itemize}

\section{3. 社会范畴（Social）：运营韧性与共享价值}

公司坚持“真情服务”底线，持续提升安全保障能力与服务品质，注重员工权益保护，积极投身社会公益事业。

\subsection{3.1 运营绩效与市场表现}

2024年，得益于旅游市场的复苏及航线网络的优化，公司核心业务指标稳步增长。

\begin{table}[H]
    \centering
    \caption{2024年度航空业务及商业运营数据}
    \label{tab:ops_stats}
    \renewcommand{\arraystretch}{1.2}
    \begin{tabular}{p{8cm} p{5cm}}
        \toprule
        \textbf{指标项目} & \textbf{统计数据} \\
        \midrule
        旗下控股及管理输出机场数量 & 9 家 \\
        总体起降架次 & 16.11 万架次 \\
        旅客吞吐量 & 2,462.62 万人次 \\
        核心门户机场行业排名 & 重回“2000万级”梯队（全国第29位） \\
        重点区域机场航线拓展 & 开通首条国际客运航线（吉隆坡） \\
        \bottomrule
    \end{tabular}
\end{table}

\subsection{3.2 安全管理与应急响应}

公司坚持“安全第一”原则，持续完善安全管理体系（SMS），提升极端天气应对能力。

\begin{itemize}
    \item \textbf{重大气象灾害应对（台风“摩羯”）}：
    \begin{itemize}
        \item \textbf{预警响应}：及时启动“防汛防风专项应急预案”，运管委统筹实施航班动态调整及航空器系留。
        \item \textbf{恢复效率}：台风过境后迅速开展适航恢复工作，9月3日至7日期间核心机场累计保障航班 1,300 架次，确保了自贸港空中通道的畅通。
    \end{itemize}
    
    \item \textbf{安全数字化管理}：
    \begin{itemize}
        \item 上线“安全智控”平台，实现安全管理流程的标准化与数字化。
        \item 平台关联制度手册 9,626 项，全年流转电子检查单 31 万份。
    \end{itemize}

    \item \textbf{安全绩效指标}：
    \begin{itemize}
        \item 一般征候万架次率：0
        \item 责任事故发生率：0
        \item 安全培训覆盖率：100\%
    \end{itemize}
\end{itemize}

\subsection{3.3 客户服务与创新}

\begin{itemize}
    \item \textbf{流程优化}：实施“一机双屏”海关安检融合查验模式，实现旅客“一次过检”，显著提升通关效率。
    \item \textbf{支付便利化}：在核心机场启用“一站式综合服务中心”，提供外卡支付绑定、通讯服务及文旅咨询，优化入境旅客体验。
\end{itemize}

\subsection{3.4 人力资源与社会贡献}

公司秉持多元包容的雇主品牌理念，并积极回馈社区。

\begin{table}[H]
    \centering
    \caption{2024年度员工与社会贡献数据}
    \label{tab:social_data}
    \renewcommand{\arraystretch}{1.2}
    \begin{tabular}{p{7cm} p{6cm}}
        \toprule
        \textbf{指标类别} & \textbf{具体数据} \\
        \midrule
        \textbf{员工构成} & \\
        员工总数 & 10,611 人 \\
        女性员工占比 & 38.9\% (4,133人) \\
        少数民族员工 & 179 人 \\
        \midrule
        \textbf{培训与发展} & \\
        年度培训总投入 & 1,119.54 万元 \\
        人均培训时长 & 105.88 小时 \\
        \midrule
        \textbf{社会投入} & \\
        乡村振兴投入 & 190 万元 \\
        志愿服务参与 & 10,000+ 人次 \\
        志愿服务总时长 & 500,000 小时 \\
        \bottomrule
    \end{tabular}
\end{table}

\section{4. 管治范畴（Governance）：合规与风险控制}

公司致力于维持高水平的企业管治标准，通过完善的治理架构与监督体系，保障公司稳健运营。

\subsection{4.1 治理架构与投资者关系}

\begin{itemize}
    \item \textbf{董事会多元化}：董事会由9名董事组成，其中独立董事3名；女性董事占比 22\%。
    \item \textbf{决策机制}：全年召开股东大会5次，董事会13次。
    \item \textbf{信息披露}：全年发布公告70份，确保信息披露的真实、准确、完整。
\end{itemize}

\subsection{4.2 内部控制与商业道德}

公司构建了“三全三化”大监督体系，强化反腐败与供应链合规管理。

\begin{table}[H]
    \centering
    \caption{2024年度合规管理绩效}
    \label{tab:gov_compliance}
    \renewcommand{\arraystretch}{1.2}
    \begin{tabular}{p{9cm} p{4cm}}
        \toprule
        \textbf{合规指标} & \textbf{执行情况} \\
        \midrule
        风险点梳理与管控 & 358 个 \\
        反腐败培训覆盖率（管理层） & 95\% \\
        反腐败培训覆盖率（员工） & 90\% \\
        反腐败培训总时长 & 400 小时 \\
        供应商廉洁承诺书签署率 & 100\% \\
        工程类供应商遴选数量 & 23 家 \\
        \bottomrule
    \end{tabular}
\end{table}

\subsection{4.3 经济绩效延伸}

在主营业务之外，公司积极参与离岛免税商业布局。2024年，公司间接参与的5家离岛免税店实现线下销售额约 \textbf{52 亿元}，市场份额占比约 \textbf{17\%}，有效促进了区域消费经济的发展。

\end{document}